\documentclass[a4paper,11pt]{article}

% 1. 中文支持 (必须使用 XeLaTeX 编译)
\PassOptionsToPackage{quiet}{fontspec}
\PassOptionsToPackage{usenames,dvipsnames}{xcolor}
\usepackage[UTF8]{ctex}

% 2. 绘图包
\usepackage{tikz}
\usetikzlibrary{shapes, backgrounds}

% 3. 布局与格式包
\usepackage{latexsym}
\usepackage[empty]{fullpage}
\usepackage[explicit]{titlesec}
\usepackage{marvosym}
\usepackage{verbatim}
\usepackage{enumitem}
\usepackage[hidelinks]{hyperref}
\usepackage{fancyhdr}
\usepackage{tabularx}
\usepackage{xcolor}

%-------------------------------------------
% 颜色定义
\definecolor{MyDarkBlue}{RGB}{0, 50, 100}

% 页边距设置
\pagestyle{fancy}
\fancyhf{}
\fancyfoot{}
\renewcommand{\headrulewidth}{0pt}
\renewcommand{\footrulewidth}{0pt}

\addtolength{\oddsidemargin}{-0.6in}
\addtolength{\evensidemargin}{-0.6in}
\addtolength{\textwidth}{1.2in}
\addtolength{\topmargin}{-.7in}
\addtolength{\textheight}{1.4in}

\urlstyle{same}
\raggedbottom
\raggedright
\setlength{\tabcolsep}{0in}

%-------------------------------------------
% TikZ 标题样式
\titleformat{\section}
  {}
  {}
  {0pt}
  {%
    \noindent
    \begin{tikzpicture}[baseline=(t.base)]
    \vspace{-2pt}
      \node[
        inner xsep=10pt,
        inner ysep=3pt,
        text=white,
        font=\large\bfseries,
        align=left,
        anchor=base west
      ] (t) at (0,0) {#1};
      \begin{scope}[on background layer]
         \fill[MyDarkBlue]
         (t.north west) --
         (t.north east) --
         ([xshift=8pt]t.south east) --
         (t.south west) -- cycle;
         \draw[MyDarkBlue, line width=2pt]
         ([xshift=-68.30pt]t.south east) --
         (\linewidth, 0 |- t.south east);
      \end{scope}
    \end{tikzpicture}%
  }
  []

\titlespacing*{\section}{0pt}{12pt}{6pt}

%-------------------------------------------
% 自定义命令

\newcommand{\resumeItem}[1]{
  \item\small{{#1 \vspace{-1pt}}}
}

\newcommand{\resumeEducationHeading}[4]{
  \vspace{-2pt}\item
    \begin{tabular*}{0.97\textwidth}[t]{l@{\extracolsep{\fill}}r}
      \textbf{#1} & #2 \\
      \textit{\small #3} & \textit{\small #4} \\
    \end{tabular*}\vspace{-5pt}
}

\newcommand{\resumeResearchHeading}[2]{
  \vspace{-2pt}\item
    \begin{tabular*}{0.97\textwidth}[t]{l@{\extracolsep{\fill}}r}
      \textbf{#1} & \textbf{#2} \\
    \end{tabular*}\vspace{-7pt}
}

\newcommand{\resumeProjectHeading}[2]{
    \item
    \begin{tabular*}{0.97\textwidth}{l@{\extracolsep{\fill}}r}
      \small#1 & #2 \\
    \end{tabular*}\vspace{-7pt}
}

\newcommand{\resumeSubHeadingListStart}{
\begin{itemize}[leftmargin=0.15in, label={}]
}

\newcommand{\resumeSubHeadingListEnd}{
\end{itemize}
}

\newcommand{\resumeItemListStart}{
\begin{itemize}[leftmargin=*, label={\small\textbullet}]
}

\newcommand{\resumeItemListEnd}{
\end{itemize}\vspace{-2pt}
}

%-------------------------------------------
%%%%%% 简历正文 %%%%%%%%%%%%%%%%%%%%%%%%%%%%

\begin{document}

%-------------------------------------------
% 个人信息

\begin{center}
    \textbf{\Huge {\ziju{0.2} 郜一贤}} \\ \vspace{8pt}

    \small
    \href{mailto:yisrealbonaparte@outlook.com}
    {\underline{yisrealbonaparte@outlook.com}}
    $|$
    +86 199-3494-2785
    $|$
    杭州电子科技大学 $\cdot$ 传播学
    $|$
    求职意向：新媒体内容相关岗位
\end{center}

%-------------------------------------------
% 教育背景

\section{教育背景}

\resumeSubHeadingListStart

\resumeEducationHeading
{杭州电子科技大学}{2023.09 -- 2027.07}
{本科 \ \textbar \ 传播学}{杭州，浙江}

\resumeItemListStart
    \resumeItem{
        \textbf{GPA: 4.5 / 5.0}
        \ \ $|\ \ $
        \textbf{主要荣誉：}
        优秀学生干部，多次获得校二、三等奖学金
    }

    \resumeItem{
        \textbf{主修课程：}
        传播学原理、非线性编辑、数据新闻等
    }

    \resumeItem{
        \textbf{英语水平：}
        CET-4 \ \textbar \ 雅思 6.0，
        具备英文文献阅读能力
    }
\resumeItemListEnd

\resumeSubHeadingListEnd

%-------------------------------------------
% 个人技能

\section{个人技能}

\begin{itemize}[leftmargin=0.15in, label={}]

\small{\item{

\textbf{设计与剪辑工具}：
\textbf{Adobe After Effects}、\textbf{Photoshop}、\textbf{达芬奇}、剪映、秀米、Canva，
熟悉视频剪辑、平面设计与排版
\\

\textbf{AI 创作工具}：
\textbf{ChatGPT}、Gemini、DeepSeek、豆包、千问，
擅长文本生成与 Prompt 设计；
\textbf{即梦}、可灵、Midjourney，
能够用于图像 / 视频内容生成
\\

\textbf{办公软件}：
WPS Office（Word / Excel / PPT）、
Audacity 音频处理
\\

\textbf{新媒体运营}：
公众号内容策划与排版、小红书账号运营、
短视频脚本撰写与剪辑、
视觉风格统一与流程标准化
\\

\textbf{个人账号成果}：
独立运营小红书账号，
产出 3 篇\textbf{点赞量破万笔记}，
\textbf{单篇最高曝光超 120 万}

}}

\end{itemize}

%-------------------------------------------
% 校园经历

\section{校园经历}

\resumeSubHeadingListStart

\resumeResearchHeading
{杭州电子科技大学校辩论协会 \ \textbar \ 宣传部部长}
{2023.09 -- 2025.09}

\resumeItemListStart

\resumeItem{
\textbf{内容运营：}
主导“HDU 辩协”官方公众号赛事内容策划与运营，
围绕赛前预告、选手专访、赛后复盘构建完整内容链路，
公众号累计覆盖用户 \textbf{805 人}，
\textbf{推文平均阅读量达 1500+}。
}

\resumeItem{
\textbf{流程优化：}
优化“选题—审核—发布—复盘”标准化流程，
统一视觉排版规范，
推动公众号\textbf{粉丝增长至 800+}。
}

\resumeItemListEnd

%-------------------------------------------

\resumeResearchHeading
{杭州电子科技大学红的家园网络工作室 \ \textbar \ 影视部副部长}
{2023.09 -- 2025.09}

\resumeItemListStart

\resumeItem{
\textbf{素材支持：}
为学校官方微博提供照片及视频素材；
通过制定标准化拍摄排期表，
统筹素材采集与后期制作，
优化跨部门协作机制，
使\textbf{素材交付效率提升 40\%}，
保障内容稳定输出。
}

\resumeItemListEnd

\resumeSubHeadingListEnd

%-------------------------------------------
% 科研竞赛

\section{科研竞赛}

\resumeSubHeadingListStart

\resumeResearchHeading
{《前沿科学》第 13 卷第 6 期 \ \textbar \ 第一作者}
{2025}

\resumeItemListStart

\resumeItem{
以\textbf{第一作者}身份在《前沿科学》（CN11-5568/N）发表论文《AI “魔改”经典影视剧的法律风险》，
主要研究 AI 生成内容在影视改编中的侵权风险、法律边界及应对策略。
}

\resumeItemListEnd

%-------------------------------------------

\resumeResearchHeading
{品牌策划赛 \ \textbar \ 团队 IP 设计成员}
{2025.03 -- 2025.05}

\resumeItemListStart

\resumeItem{
负责民宿品牌 \textbf{“灵溪山居”}
的 IP 形象与文创产品设计，
融合山水意象与年轻化叙事，
原创 IP 角色 \textbf{“溪溪”}
及配套视觉体系。
}

\resumeItem{
主要借助即梦、Midjourney 等 AI 工具完成视觉生成，
并结合 ChatGPT 进行 Prompt 设计与内容优化。
}

\resumeItem{
围绕文创产品、短视频矩阵与沉浸式体验馆
构建多维传播方案，
项目最终获
\textbf{\color{MyDarkBlue}国家级二等奖}。
}

\resumeItemListEnd

%-------------------------------------------

\resumeResearchHeading
{全国大学生广告艺术大赛 \ \textbar \ 剪辑 / 演员}
{2025.03 -- 2025.06}

\resumeItemListStart

\resumeItem{
参与大广赛营销创客微短剧赛道，
作品《UU 跑腿迷案》
获 \textbf{\color{MyDarkBlue}国赛三等奖}。
}

\resumeItem{
担任剪辑与演员，
部分画面使用\textbf{即梦 AI 生成}，
部分音频采用 AI 配音；
通过剪辑节奏控制与表演呈现品牌“高效应急”价值，
完整参与团队创作流程。
}

\resumeItemListEnd

%-------------------------------------------

\resumeResearchHeading
{EOVS 工创赛 \ \textbar \ 团队制表成员}
{2024.09 -- 2024.11}

\resumeItemListStart

\resumeItem{
作为企业运营虚拟仿真赛项核心成员，
主导财务与市场数据制表分析，
优化季度经营报表框架，
\textbf{支撑团队完成 8 个季度运营策略迭代}，
最终获
\textbf{\color{MyDarkBlue}省赛三等奖}。
}

\resumeItemListEnd

\resumeSubHeadingListEnd

%-------------------------------------------
% 实习与项目经历

\section{实习与项目经历}

\resumeSubHeadingListStart

\resumeResearchHeading
{尚岸（杭州）教育信息技术有限责任公司 \ \textbar \ 设计部实习生}
{2025.07 -- 2025.08}

\resumeItemListStart

\resumeItem{
负责\textbf{人教版小学、初中教材交互板块音频导入}，
使用 \textbf{Audacity} 分割音频文件并编号入库，
确保音频与教材内容精准匹配，
高效完成交付任务并保障交付质量，实习表现并获得公司认可。
}

\resumeItem{
同时使用 \textbf{Photoshop} 设计小学英语单词卡片，
输出标准化教学辅助素材。
}

\resumeItemListEnd

%-------------------------------------------

\resumeResearchHeading
{浙江大学出版社数字教材项目 \ \textbar \ 编辑}
{2025.03 -- 2025.07}

\resumeItemListStart

\resumeItem{
负责在浙大数字教材网站录入书目信息，
逐项核对教材内容与格式，
\textbf{累计录入文史、理工、经管等学科书目 10 余本}，
\textbf{录入准确率达 95\% 以上}。
}

\resumeItemListEnd

\resumeSubHeadingListEnd

%-------------------------------------------
% 荣誉奖项

\section{荣誉奖项}

\resumeSubHeadingListStart

\resumeProjectHeading
{\textbf{品牌规划赛} 国家级二等奖（团队 IP 设计成员）}
{2025}

\resumeProjectHeading
{\textbf{全国大学生广告艺术大赛} 国赛三等奖（剪辑 / 演员）}
{2025}

\resumeProjectHeading
{\textbf{EOVS 工创赛} 省赛三等奖（团队制表成员）}
{2024}

\resumeProjectHeading
{\textbf{校二、三等奖学金}（多次获得）
\ \& \
\textbf{优秀学生干部}}
{本科}

\resumeProjectHeading
{《前沿科学》第 13 卷第 6 期
\textbf{第一作者}发表论文}
{2025}

\resumeSubHeadingListEnd

%-------------------------------------------

\end{document}